% Transcription — fonds Alexandre Grothendieck, shelfmark 161-2, batch 6 (pages 101-111).
% Established from the facsimile archives/batches/161-2/batch-06.pdf (a mirror of
% https://grothendieck.umontpellier.fr/161-2.pdf), per the skill
% transcribe-grothendieck. The folder's last eleven pages. Pages 107 and 110
% are blank and absent. Three runs: pp. 101-106, the end of the fair copy he
% paginated ⑤-⑩ (①-④ are pp. 97-100 of batch 5) — the questions a')-d') on
% λ-generation, λ-variable (η_{λλ'}), duality, and §8, structures defined by
% projective limits: the 2-cartesian squares (8.2.3), (8.2.5), (8.2.10), the
% functor α : T(Ĉ) → Hom(C°, S_T), R° ⊂ S ⊂ \hat{R°}, and a summary that
% stops mid-sentence; pp. 108-109, a faint pencil sheet on the classifying
% topos Ĉ = B_T, the universal structure O_B, f^*(O_B) and the pull-back of a
% stack, with T(Ens) = Ind(C°) and Pro(C) ≃ T(Ens)°; p. 111, a squared sheet
% headed « Topos modulaire pour des structures algébriques sur objets d'un
% topos fixé X », the 2-fibre product B ×_{B_I} X. The pencil pages and p. 111
% carry many \ill{}: the formulas pass, the connective prose often does not.
% Pass: Fable 5.1 (claude-fable-5-1), 2026-09-04 — first pass, unchecked against the pages by a human.
\documentclass[11pt,a4paper]{article}
% Le préambule commun à toutes les transcriptions du fonds Grothendieck.
%
% Il tient en une idée : **l'incertitude se compose, elle ne se résout pas.**
% Une transcription qui lisse un mot douteux a détruit la seule chose pour
% laquelle on la faisait — un lecteur doit pouvoir savoir, à chaque endroit,
% ce qui était sur le feuillet et ce qui a été ajouté. Les six macros qui
% suivent sont donc l'appareil critique, et non de la décoration.
%
% Les mêmes commandes sont reconnues par `scripts/render.mjs`, qui produit la
% vue de lecture du site. Une macro ajoutée ici doit l'être aussi là-bas,
% faute de quoi le rendu échouera bruyamment — ce qui est voulu : un rendu
% silencieusement faux est pire qu'un rendu absent.

\ProvidesPackage{grothendieck}[2026/08/08 Transcriptions du fonds Grothendieck]

\usepackage{amsmath,amssymb,amsthm}
\usepackage{mathrsfs}
\usepackage{stmaryrd}
% Les diagrammes commutatifs sont partout dans ce fonds : c'est la langue
% même de la géométrie algébrique de Grothendieck, et une transcription qui
% les rendrait en prose ne transcrirait rien.
\usepackage{tikz-cd}
% Français par défaut — c'est la langue des feuillets — mais l'anglais est
% chargé aussi : la traduction et le résumé sont des documents anglais, et la
% typographie française (espaces avant les deux-points) y serait une faute.
% Ces documents basculent par \selectlanguage{english} en tête de corps.
\usepackage[english,french]{babel}
\usepackage{fontspec}
\usepackage{geometry}
\geometry{a4paper,margin=25mm,marginparwidth=28mm}
\usepackage{marginnote}
\usepackage{xcolor}
% ulem, not soul. `soul`'s \st and \ul are built on a character-by-character
% scan that XeLaTeX's font machinery breaks: the document fails with an
% undefined control sequence rather than degrading. ulem does the same two jobs
% and is Unicode-engine safe. `normalem` stops it hijacking \emph, which the
% transcriptions use for Grothendieck's own emphasis.
\usepackage[normalem]{ulem}
\usepackage{hyperref}
\hypersetup{colorlinks=true,linkcolor=black,urlcolor=black,pdfborder={0 0 0}}

% --- Métadonnées du lot -----------------------------------------------------
% Reprises telles quelles par le script de rendu, qui les affiche en tête de
% la vue de lecture. Elles fixent ce que le fichier prétend couvrir : sans
% elles, une transcription flotte sans point d'attache dans un fonds de seize
% mille feuillets.

\newcommand{\folder}[1]{\gdef\grothFolder{#1}}
\newcommand{\batch}[1]{\gdef\grothBatch{#1}}
\newcommand{\leaves}[2]{\gdef\grothFirst{#1}\gdef\grothLast{#2}}
\newcommand{\foldertitle}[1]{\gdef\grothFoldertitle{#1}}
\gdef\grothFolder{?}\gdef\grothBatch{?}\gdef\grothFirst{?}\gdef\grothLast{?}\gdef\grothFoldertitle{}

% --- Le résumé --------------------------------------------------------------
%
% Il ouvre la lecture modernisée, sous le titre « Résumé », et il est composé
% en petit. Ce n'est pas un ornement : c'est le seul passage du document qui
% ne soit pas des mathématiques, et un lecteur doit pouvoir le reconnaître
% avant de l'avoir lu — puis le sauter s'il connaît déjà le sujet, ou s'y
% arrêter s'il ne le connaît pas. Le filet qui le suit marque la reprise du
% corps.

\newenvironment{resume}
  {\small\setlength{\parskip}{0.45em}}
  {\par\medskip\hrule\medskip}

% --- Mots-clés ---------------------------------------------------------------
%
% En anglais, en clôture du résumé de la lecture modernisée : le vocabulaire
% moderne sous lequel chercher ce que les feuillets construisent. Le manifeste
% du site (`npm run manifest`) les extrait pour étiqueter la cote entière —
% cette ligne est la source unique des tags, il n'y a pas de fichier à côté.
\newcommand{\keywords}[1]{\par\smallskip\noindent
  {\footnotesize\textsc{Keywords} --- \textit{#1}}\par}

% --- L'appareil critique ----------------------------------------------------

% Le numéro de feuillet, en marge. C'est l'ancre qui permet de régler un
% désaccord sur une formule en regardant *un* feuillet plutôt qu'une chemise
% de six cents. Le site s'en sert aussi pour tourner les pages du fac-similé
% quand on fait défiler la transcription.
\newcommand{\leaf}[1]{%
  \marginnote{\footnotesize\color{gray!70}\textsf{#1}}%
  \label{leaf:#1}\ignorespaces}

% Illisible : on ne devine pas. Un mot inventé qui se lit comme les autres est
% le pire résultat possible.
\newcommand{\ill}{\textcolor{red!60!black}{[\,\ldots\,]}}

% Lecture douteuse : le mot est donné, et signalé comme incertain.
\newcommand{\uncertain}[1]{\uline{#1}}

% Ajout de l'éditeur — une lettre manquante, un mot sous-entendu.
\newcommand{\add}[1]{\textcolor{blue!50!black}{[#1]}}

% Biffé par Grothendieck lui-même. Ce qu'il a rayé fait partie du document :
% une rature dit souvent où la pensée a changé de direction.
\newcommand{\struck}[1]{\sout{#1}}

% Note du transcripteur, sur l'état matériel du feuillet.
\newcommand{\note}[1]{\par\smallskip{\small\color{gray!60!black}\textit{[#1]}}\par\smallskip}

% Note marginale de Grothendieck, distincte du corps du texte.
\newcommand{\marginal}[1]{\par\smallskip\noindent\hspace{1em}%
  {\small\color{gray!60!black}#1}\par\smallskip}

% --- Titre ------------------------------------------------------------------

\AtBeginDocument{%
  \begin{center}
    {\large\bfseries Fonds Alexandre Grothendieck --- cote n\textsuperscript{o}~\grothFolder}\\[2pt]
    {\itshape\grothFoldertitle}\\[4pt]
    {\small feuillets \grothFirst--\grothLast\ (lot \grothBatch)}\\[2pt]
    {\footnotesize\color{gray!60!black}Transcription établie d'après le fac-similé numérisé
     par l'Université de Montpellier.\\
     Les lectures incertaines sont soulignées, les illisibles notées [\,\ldots\,],
     les ajouts entre crochets.}
  \end{center}
  \vspace{1em}}

\endinput


\folder{161-2}
\batch{6}
\pages{101}{111}
\dating{[vers 1963-1973]}
\watermark{Édition de démonstration}
\foldertitle{Algèbre universelle [ou catégories] : notes manuscrites (s.d.)}

\begin{document}

\page{101}
\note{author's p. ⑤ ; la copie paginée ①--④ du lot précédent (pp. 97--100)
se poursuit ici jusqu'à ⑩ (p. 106) ; les alinéas 6, 7, 8 sont numérotés
dans des carrés}

\underline{5.4.}\ L'essentiel \struck{\ill{}} se décompose en les questions

\begin{itemize}
  \item[a')] $\tau^{I}$ est-il représentable par une catégorie $R_{I}$
    (la $\lambda$-catégorie libre engendrée par $I$ objets) qui soit
    $\lambda$-engendrée (au sens de (5.3)) par
    \uncertain{$\mathrm{Im}(I \to \mathrm{Ob}\,R_{I})$} ?
  \item[b')] Soient $R \in \mathrm{Ob}\,(\mathrm{Cat}_{\lambda})$ \note{un mot biffé entre $\mathrm{Ob}$ et la parenthèse},
    $A \in \mathrm{Fl}\,R$. Y a-t-il une $\lambda$-catégorie $R_{A}$
    \add{$\in \mathrm{Ob}\,\mathrm{Cat}_{\lambda}$} « universelle » pour
    rendre inversibles les flèches $\in A$, et $R_{A}$ est-il
    $\lambda$-engendré par \uncertain{$\mathrm{Im}\,R$} ?
    \marginal{cf. note \uncertain{avant} b)}
  \item[\uncertain{c'})] \note{l'encre se lit plutôt « b'' »} Soit
    $R \in \mathrm{Ob}(\mathrm{Cat}_{\lambda})$, $B$ un \struck{\ill{}}
    petit ensemble, et $\alpha : B \to \mathrm{Ob}\,R
    \times \mathrm{Ob}\,R$. \struck{\ill{}} Y a-t-il une $\lambda$-cat.
    $R' = R_{B,\alpha} \in \mathrm{Ob}\,\mathrm{Cat}_{\lambda}$
    universelle \struck{\ill{}} munie d'un $\lambda$-foncteur
    $R \xrightarrow{\ i\ } R'$ et d'applications $B \to
    \mathrm{Fl}\,R'$ \note{un mot biffé devant $\mathrm{Fl}\,R'$} telle que le composé avec $\mathrm{Fl}\,R' \to
    \mathrm{Ob}\,R' \times \mathrm{Ob}\,R'$ soit le composé de $\alpha$
    avec $\mathrm{Ob}\,i \times \mathrm{Ob}\,i$. De plus, $R'$ est-il
    $\lambda$-engendré par \uncertain{$\mathrm{Im}\,i$} ?
  \item[d')] Soient $(R_{j})_{j \in J}$ une petite famille d'objets de
    $(\mathrm{Cat}_{\lambda})$, munies d'applications $I
    \xrightarrow{\ \varphi_{j}\ } \mathrm{Ob}\,R_{j}$. Y a-t-il une
    solution au pb universel d'envoyer les $R_{j}
    \xrightarrow{\ \psi_{j}\ } R$, $R \in \mathrm{Ob}\,\mathrm{Cat}_{\lambda}$,
    avec \struck{des applications} \add{un} \uncertain{syst. transitif d'}
    isomorphismes entre les \struck{\ill{}} $\psi_{j} \circ \varphi_{j}$
    ($I$ considéré comme cat. discrète\ldots), et $R$ est-il
    $\lambda$-engendré par la réunion des images des $\psi_{j}$ ?
\end{itemize}

\ill{} \uncertain{envie} des preuves \underline{5.5.}\ Soit $S \subset
\mathrm{Ob}\,R$, $R \in \mathrm{Ob}\,C_{\lambda}$ \note{$C_{\lambda}$ ici
et deux fois plus bas, abréviation manifeste de
$\mathrm{Cat}_{\lambda}$}, supposons que $\forall R' \in \mathrm{Ob}\,
C_{\lambda}$, et $F, G \in \mathrm{Hom}_{\lambda}(R, R')$,
\[
  \mathrm{Hom}(F, G) \longrightarrow \prod_{s \in S}
  \mathrm{Hom}(F(s), G(s))
\]
est injectif. Alors $R$ est-il $\lambda$-engendré par $S$ ? (la
réciproque étant évidemment vraie\ldots)

\page{102}
\note{author's p. ⑥}

\underline{5.6.}\ Partons de $R \in \mathrm{Ob}\,\mathrm{Cat}_{\lambda}$,
avec $I \xrightarrow{\ \varphi\ } \mathrm{Ob}\,R$ tels que $R$ soit
$\lambda$-engendré par $I$. Considérons la cat. 2-\struck{\ill{}}fibrée
$T$ sur $(\mathrm{Cat}_{\lambda})$ définie par $R$, et \ill{}
\uncertain{supplémentaire} $T \xrightarrow{\ \varphi\ } \tau^{I}$ déduit
de $\varphi$. \uncertain{Prouver}, que $(T, \varphi)$ est un
$\lambda$-$I$-type ? (Il faut peut-être renforcer la notion de
sous-catégorie $\lambda$-engendrée, en tenant compte des flèches, pas
seulement des objets\ldots ?)

\underline{5.7.}\ Si on laisse tomber $I$, une \struck{\ill{}} catégorie
2-cofibrée $T$ sur $(\mathrm{Cat}_{\lambda})$ est un $\lambda$-type, si
c'est un $\lambda$-$I$-type pour \add{$I$ et} $\varphi : T \to \tau^{I}$
convenables. Ceci signifie donc (si \ill{} \uncertain{ce qui précède}
\ill{}) que $T$ est $\lambda$-représentable par un $R \in
\mathrm{Ob}(\mathrm{Cat}_{\lambda})$, \uncertain{i.e.}
$\mathrm{Ob}\,R$ admet une petite partie qui soit $\lambda$-génératrice
(en un sens convenable à dégager).
\marginal{Question \ill{} \uncertain{diagrammes finis}, dans un $R \in
\mathrm{Ob}\,\mathrm{Cat}_{\lambda}$ \ill{} d'un $I \to \mathrm{Ob}\,R$
\ill{} \uncertain{type} \ill{} catégorie ?}
\note{la note marginale est écrite en diagonale, sur une dizaine de
lignes, et reste en grande partie illisible ; seuls les symboles
passent}

\textbf{[6]} \underline{$\lambda$-variable}. On devrait trouver (pour
$\lambda \subset \lambda'$) un 2-foncteur
\begin{equation}
  (\lambda\text{-types}) \xrightarrow{\ \eta_{\lambda\lambda'}\ }
  (\lambda'\text{-types})
  \tag{6.1}
\end{equation}
\note{à droite, entre parenthèses : $\eta_{\lambda,\lambda'}(T_{\lambda})(C)
= T_{\lambda}(C)$ si $C \in \mathrm{Ob}\,\mathrm{Cat}_{\lambda'}$
($\subset \mathrm{Ob}\,\mathrm{Cat}_{\lambda}$)}
qui correspondent à un 2-foncteur
\begin{equation}
  (\mathrm{Cat}_{\lambda}) \xrightarrow{\ \eta_{\lambda\lambda'}\ }
  (\mathrm{Cat}_{\lambda'})
  \tag{6.2}
\end{equation}
\note{le symbole sur la flèche de (6.2) porte peut-être un accent}
« $\lambda'$-catégories \uncertain{analysantes} ». Il faudrait donner
des conditions moyennant lesquelles (6.1) est pl. fidèle. [\underline{Ex}
\struck{\ill{}} $\varprojlim$ finies, $\lambda' = \varprojlim$ quelc.
($\bar{\eta}(C) = \mathrm{Pro}(C)$ ?) ; $\lambda = \varinjlim$ finies,
$\lambda' = \varinjlim$ quelc. ($\bar{\eta}(C) = \mathrm{Ind}(C)$ ?)]

\note{en marge : « ! \textbar\textbar »} $\eta_{\lambda,\lambda'}$
est-il \struck{\ill{}} pl.\ fidèle, i.e.\ $\mathrm{Hom}_{\lambda}(T_{1},
T_{2}) \to \mathrm{Hom}_{\lambda'}(T_{1}, T_{2})$ est-il \ill{} une
équivalence pour $T_{1}$ et $T_{2}$ des $\lambda$-types ? Sans doute
faux, mais \ill{} pl.\ fidèle si $\overrightarrow{\lambda} = \emptyset$.

\textbf{[7]} \underline{Dualité} : \underline{Opposé} d'un « type »
$\lambda = (\overleftarrow{\lambda}, \overrightarrow{\lambda})$ : on
échange les types de diagrammes utilisés pour les $\varprojlim$
\add{et} $\varinjlim$. Alors pour tt $\lambda$-type $T_{\lambda}$, on
définit un $\lambda^{\circ}$-type $T^{\circ}_{\lambda}$ \ill{}
\begin{equation}
  T^{\circ}_{\lambda}(C) \cong T_{\lambda}(C^{\circ})^{\circ}
  \tag{7.1}
\end{equation}
\note{les indices $\lambda$ de (7.1), ici et en haut de la page suivante,
sont surchargés}

\page{103}
\note{author's p. ⑦}
i.e.\ on a une équivalence de Cat. 2-cofibrées
\begin{equation}
  T^{\circ}_{\lambda} \cong (T_{\lambda})^{\circ} \circ (\ldots)
  \tag{7.1}
\end{equation}
\note{le dernier facteur, entre parenthèses, est indistinct ; on lit
peut-être « pass. à l'opp. » ; au-dessous, un numéro (7.2) biffé}
On trouve ainsi une 2-équivalence
\begin{equation}
  (\lambda\text{-types}) \cong (\lambda^{\circ}\text{-types})
  \tag{7.3}
\end{equation}
(et kif kif avec $I$
\begin{equation}
  (\lambda\text{-}I\text{-types}) \cong
  (\lambda^{\circ}\text{-}I\text{-types})\ )
  \tag{7.4}
\end{equation}
compatible, via les cat. représentatives, --
\begin{equation}
  (\mathrm{Cat}_{\lambda}) \cong (\mathrm{Cat}_{\lambda^{\circ}}),
  \qquad C \longmapsto C^{\circ}
  \tag{7.6}
\end{equation}
\note{le numéro est surchargé, 5 ou 6}

\textbf{[8]} Cas où $\overrightarrow{\lambda} = \emptyset$ (structures
définies via $\varprojlim$). $T$ désignera une $\lambda$-structure.

\underline{8.1.}\ Pour tout $C \in \mathrm{Ob}\,\mathrm{Cat}_{\lambda}$,
les \struck{limites} $\varprojlim$ de type $\lambda$ existent dans
$T(C)$ ; plus généralement, \struck{les lim} \add{pour} tout type $d$
\struck{\ill{}} \add{de} $\varprojlim$ qui existent dans $C$ existent
aussi dans $T(C)$. [Écrire $T(C) \cong \mathrm{Hom}_{\lambda}(R_{T}, C)$
et noter que \struck{\ill{}} l'hyp.\ sur $C$ implique que
$\mathrm{Hom}(R_{T}, C)$ \note{au-dessus : « se calculant
\uncertain{argument par argument} »} \struck{\ill{}} a des
$\varprojlim$ de type $d$, et que \struck{\ill{}} $\mathrm{Hom}_{\lambda}
(R_{T}, C)$ \ill{} par \uncertain{icelles} $\varprojlim$.] De plus, pour
tt $\xi \in R_{T}$, \struck{\ill{}} \uncertain{$\Gamma_{C}$}\,$: T(C) \to
C$ commute aux $\varprojlim$ de type $d$.

\underline{8.1.1.}\ Plus généralement ($\lambda$ quelconque, sans néc.\
$\overrightarrow{\lambda} = \emptyset$). Si \add{$C \in
\mathrm{Ob}\,\mathrm{Cat}_{\lambda}$ et si} \struck{\ill{}} $d$ est un
type de diagrammes tel que
\begin{itemize}
  \item[a)] $C$ admet des $\varprojlim$ de type $d$ (p.ex.\ $d \in
    \overleftarrow{\lambda}$)
  \item[b)] dans $C$, les $\varprojlim$ de type $d$ commutent aux
    $\varinjlim$ de type $\in \overrightarrow{\lambda}$
\end{itemize}
Alors \struck{\ill{}} $T(C)$ admet des $\varprojlim$ de type $d$, et
\struck{\ill{}} icelles $T(C) \simeq \mathrm{Hom}_{\lambda}(R^{\lambda}_{T},
C)$ \struck{\ill{}} se calculent « argument par argument ». \add{Pour}
tout\add{e} \struck{\ill{}} construction $\rho$ rel.\ à $T$, le foncteur
correspondant
\[
  \rho : T(C) \longrightarrow C
\]
\marginal{de plus si $C'$ est \ill{} et \ill{} $C \to C'$ commute aux
$\varprojlim$ de type $d$, alors \struck{$T(C) \to T(C')$} \ill{}
$T(C) \to T(C')$ \ill{} ; kif-kif si $T \to T'$, $T(C) \to T'(C)$
\ill{} commute aux \ill{} de type $d$ \ill{}}
\note{note marginale en diagonale au bas de la page, en partie
illisible ; la phrase « le foncteur correspondant $\rho$ » se poursuit
en haut de la page suivante}

\page{104}
\note{author's p. ⑧}
commute aux $\varprojlim$ de type $d$. Bien sûr, on a un énoncé dual
pour les $\varinjlim$.

\underline{8.2.}\ \marginal{(remarque)} Notons que pour \struck{\ill{}}
quelconque \struck{\ill{}} \add{et un $\lambda$-type $T$,} il est vrai
que pour tt $\lambda$-foncteurs \note{en insertion, entre accolades :
« pl. fid. (fid.) les foncteurs »}
\begin{equation}
  i : C \longrightarrow C'
  \tag{8.2.1}
\end{equation}
\begin{equation}
  T(i) : T(C) \longrightarrow T(C')
  \tag{8.2.2}
\end{equation}
\struck{\ill{}} \add{est} pl.\ fidèle (fid.). Si $T$ est muni de $b :
T \to \tau^{I}$ (« base ») \underline{fidèle}, i.e.\ $R =
R^{\lambda}_{T}$ \note{au-dessus : « qui 2-représente $T$ »} est muni
de $\varphi : I \to \mathrm{Ob}\,R$ \note{l'indice de $R$ est biffé}, alors
\struck{le carré} dans le cas pl.\ fidèle

\begin{tikzcd}
  T(C) \arrow[r] \arrow[d] & T(C') \arrow[d] \\
  C^{I} \arrow[r, hook] & C'^{I}
\end{tikzcd}

\note{numéroté (8.2.3)} est 2-cartésien [Ceci signifie aussi, en termes
de $R$, que la \uncertain{ss}-catégorie pleine des
$\mathrm{Hom}_{\lambda}(R, C)$ de $\mathrm{Hom}_{\lambda}(R, C')$
\struck{s'identifie aux} \add{est formée des} foncteurs $R
\xrightarrow{\ \xi\ } C'$ tels que $\xi(\varphi(i)) \in C$ pour tt $i
\in I$ : en effet, comme $I$ $\lambda$-engendre $R$, il s'ensuit que
$\xi$ (qui commute aux $\varprojlim$ de type $\lambda$)
\uncertain{envoie} $R$ dans l'image \uncertain{essentielle} \add{de}
$C$.] \note{en insertion, sur deux lignes serrées :} En termes imagés,
\ill{} où $(E_{i})_{i \in I}$ est une famille d'objets de $C$, « il
revient au même » de se donner \marginal{dessus une structure d'espèce
$T$ sur les $E_{i}$ \ill{} et une \ill{} sur les images \ill{}}
\note{la fin de la phrase court dans la marge gauche, en partie
illisible}

Appliquant ceci, pour $C \in \mathrm{Ob}\,\mathrm{Cat}_{\lambda}$, au
plongement canonique
\begin{equation}
  i : C \hookrightarrow \hat{C}
  \tag{8.2.4}
\end{equation}
\note{à droite, souligné d'un trait ondulé : « qui est un
$\lambda$-foncteur si $\lambda = (\overleftarrow{\lambda},
\emptyset)$ ; »} on trouve le foncteur

\begin{tikzcd}
  T(C) \arrow[r] \arrow[d] & T(\hat{C}) \arrow[d] \\
  C^{I} \arrow[r] & \hat{C}^{I}
\end{tikzcd}

\note{numéroté (8.2.5) ; à droite : « (2-cartésien) », puis, entouré :
« On suppose $\overrightarrow{\lambda} = \emptyset$ »}
\struck{Or supposons ici \ill{} $\hat{C}$ \ill{} $\mathrm{Ob}\,
\mathrm{Cat}_{\lambda}$ \ill{} catégorie \ill{} i.e.\ que $\hat{C}$
\ill{} si $C$ est $\mathcal{U}$-petite \ill{} Dans $T(\hat{C})$ \ill{}
$X \in$} Considérons le foncteur contravariant canonique de $\hat{C}$
dans $\mathrm{Hom}(\hat{C}, (\mathrm{Ens}))$, \struck{\ill{}} \add{sa}
restriction à $C$ \struck{\ill{}} moyennant $i$ est un foncteur
canonique $C^{\circ} \to \mathrm{Hom}_{\lambda}(\hat{C},
(\mathrm{Ens}))$, et qui s'en va en \struck{\ill{}} fait dans la
\struck{$\hat{C}$} ss-catégorie pleine des foncteurs qui commutent
\struck{\ill{}} à tt type \add{« petit »} de lim (ind ou proj), donc en
particulier
\begin{equation}
  C^{\circ} \longrightarrow \mathrm{Hom}_{\lambda}(\hat{C},
  (\mathrm{Ens})), \qquad X \longmapsto (F \mapsto F(X))
  \tag{8.2.6}
\end{equation}
D'où on a \ill{}
\[
  \mathrm{Hom}_{\lambda}(\hat{C}, (\mathrm{Ens})) \longrightarrow
  \mathrm{Hom}(T(\hat{C}), T((\mathrm{Ens})))
\]
\note{les deux $T$ du membre de droite portent un indice biffé}
\marginal{\underline{NB} L'hyp.\ « $\hat{C}$ une
$\mathcal{U}$-catégorie » \ill{} est \uncertain{essentielle} pour la
\uncertain{représentabilité} \ill{} (8.2.\uncertain{9}) \ill{} $C$
\ill{} \uncertain{catégorie} \ill{} $T(\hat{C})$ \ill{} de
\uncertain{l'« absolu »} \ill{}}
\marginal{on suppose \ill{} des univers \ill{} $\lambda$ quelc.}
\note{deux notes marginales à gauche, la première longue et en grande
partie illisible ; leur place dans le texte n'est pas marquée}

\page{105}
\note{author's p. ⑨}
d'où en composant avec (8.2.6) un foncteur canonique
\begin{equation}
  C^{\circ} \longrightarrow \mathrm{Hom}(T(\hat{C}), T(\mathrm{Ens}))
  \tag{8.2.7}
\end{equation}
ou, de façon équivalente, un foncteur canonique
\begin{equation}
  T(\hat{C}) \xrightarrow{\ \alpha\ \approx\ } \mathrm{Hom}(C^{\circ},
  S_{T}) \qquad \text{où } S_{T} = T(\mathrm{Ens})
  \tag{8.2.8}
\end{equation}
\note{le second argument de $\mathrm{Hom}$ a d'abord été écrit
$(\mathrm{Ens})$, biffé}
Si on a un ensemble d'indices des bases pour $T$, on aura un carré
commutatif

\begin{tikzcd}[column sep=large]
  T(\hat{C}) \arrow[r, "\alpha\ \approx"] \arrow[d, "b_{\hat{C}}"'] &
  \mathrm{Hom}(C^{\circ}, S_{T})
  \arrow[d, "{\mathrm{Hom}(C^{\circ},\, b_{(\mathrm{Ens})})}"] \\
  \hat{C}^{I} \arrow[r, "\cong"'] & \mathrm{Hom}(C^{\circ},
  (\mathrm{Ens}))^{I}
\end{tikzcd}

\note{numéroté (8.2.9)}
\marginal{$\mathrm{Hom}_{\lambda}(R, \hat{C}) \cong \mathrm{Hom}(R
\times C^{\circ}, (\mathrm{Ens})) \cong \mathrm{Hom}(C^{\circ},
\mathrm{Hom}_{\lambda}(R, (\mathrm{Ens})))$}
\note{formule encadrée dans la marge gauche ; l'indice du $\mathrm{Hom}$
extérieur et son premier argument sont surchargés, et un renvoi
« (8.2.8) » est écrit sous le dernier terme}
le foncteur \struck{\ill{}} $\alpha$ est tjs une équivalence,
\struck{donc} \add{si $\overrightarrow{\lambda} = \emptyset$}
\struck{\ill{}} \add{le} carré 2-cartésien

\begin{tikzcd}[column sep=large]
  T(C) \arrow[r, hook] \arrow[d, "b_{C}"'] & \mathrm{Hom}(C^{\circ},
  S_{T}) \arrow[d, "{\varphi \mapsto (b_{i,\mathrm{Ens}} \circ
  \varphi)_{i \in I}}"] \\
  C^{I} \arrow[r, hook] & \hat{C}^{I}
\end{tikzcd}

\note{numéroté (8.2.10) ; à droite : « (2-cartésien \ill{}) »,
« ($\overrightarrow{\lambda} = \emptyset$) », le reste raturé ; dans la
marge gauche, entouré : $\overrightarrow{\lambda} = \emptyset$}

Ou, en termes imagés : \add{Se} donner une structure d'espèce $T$
\add{dans $C$} revient kif-kif \add{à la} \ill{} dans $\hat{C}$, i.e.\
à la donnée, pour $(E_{i}) \in \mathrm{Ob}\,C^{I}$, d'une
\struck{\ill{}} structure d'espèce $T$ sur le \struck{\ill{}}
\add{foncteur} variant fonctoriellement avec \ill{} \struck{\ill{}}
variant : la donnée, \struck{pour chaque} \add{fonctoriellement}
\struck{la donnée} pour tt $Z \in \mathrm{Ob}\,C$ d'une \struck{\ill{}}
\add{structure d'espèce $T$ sur} $\big(E_{i}(Z)\big)_{i \in Z} =
\big(\mathrm{Hom}(Z, E_{i})\big)_{i \in I}$ \note{sic : « $i \in Z$ »
au premier indice} \struck{\ill{}} avec $Z$.
\note{le passage est écrit sur deux colonnes avec des insertions
interlinéaires ; l'ordre de lecture est reconstitué}

Lorsque $\alpha$ est une \struck{\ill{}} équivalence, \note{en
insertion : « (pour tt $C$ tel que \ill{} $\mathcal{U}$-catégorie) »}
\struck{alors} \struck{structure} $T(C)$ est connu (à équivalence
près (\ldots)) lorsque l'on connaît
\begin{equation}
  \begin{cases}
    S = S_{T} = T(\mathrm{Ens}), \text{ avec les foncteurs} \\
    \beta_{i} : S_{T} \longrightarrow (\mathrm{Ens})
  \end{cases}
  \tag{8.2.11}
\end{equation}
\note{à droite des accolades : « $\beta = b_{\mathrm{Ens}}$ », l'indice
surchargé}
comme la sous-catégorie pleine des $\mathrm{Hom}(C^{\circ}, S_{T})$
formée des $\varphi : C^{\circ} \to S_{T}$ tels que les
\uncertain{$\beta_{i} \circ \varphi$} soient représentables.

\underline{8.2.12.}\ Si on veut une car.\ intrinsèque, on peut dire
aussi que c'est la sous-catégorie pleine des $\varphi : C^{\circ} \to
S_{T}$ tels que $\beta_{\mathrm{Ens}} \circ \varphi$
\struck{représentables} \add{représentable} pour tt $\beta \in
R^{\lambda}_{T}$. Notons que le foncteur
\begin{equation}
  \beta \longmapsto \beta_{\mathrm{Ens}} : R = R^{\lambda}_{T}
  \longrightarrow \mathrm{Hom}(S_{T}, (\mathrm{Ens})) =
  \widehat{S^{\circ}_{T}}
  \tag{8.2.13}
\end{equation}
\note{sous le premier argument de $\mathrm{Hom}$, relié par un double
trait : « $S_{T}$ » ; le membre de droite est suivi d'un numéro
surchargé, (8.2.15) ou (8.2.13)}
est aussi le foncteur évident
\begin{equation}
  R \longrightarrow \mathrm{Hom}(\mathrm{Hom}_{\lambda}(R, \mathrm{Ens}),
  (\mathrm{Ens}))
  \tag{8.2.14}
\end{equation}
Or, si on considère $S = \mathrm{Hom}_{\lambda}(R, (\mathrm{Ens}))$
comme ss-catégorie pleine de $\widehat{R^{\circ}}$, \underline{celle-ci
contient $R^{\circ}$} (\add{identifié à} la cat.\ des foncteurs
représentables de $R$ dans $(\mathrm{Ens})$) donc \ldots\
\struck{\ill{}} $\widehat{R^{\circ}}$,

\page{106}
\note{author's p. ⑩}
\begin{equation}
  R^{\circ} \subset S \subset \widehat{R^{\circ}}
  \tag{8.2.15}
\end{equation}
\struck{\ill{}} et le foncteur
\[
  \pi : R \times S \longrightarrow (\mathrm{Ens})
\]
est donné par
\begin{equation}
  \pi(c, \xi) = \xi(c) =
  \mathrm{Hom}_{\widehat{R^{\circ}}}(c, \xi) = \mathrm{Hom}_{S}(c, \xi)
  \tag{8.2.16}
\end{equation}
i.e.\ \struck{\ill{}} \add{(8.2.13) est} \note{au-dessous, en petit et
biffé : « l'opposé »} des foncteurs \struck{foncteur} \add{le}
composé des deux foncteurs pleinement fidèles
\begin{equation}
  R \hookrightarrow S^{\circ}
  \hookrightarrow \widehat{S^{\circ}} =
  \mathrm{Hom}(S, (\mathrm{Ens}))
  \tag{8.2.17}
\end{equation}
\note{après le signe $=$ de (8.2.16), un terme biffé ; dans (8.2.17) le premier terme porte un exposant $\circ$ surchargé, sans doute
biffé ; sous la première flèche : « (8.2.15) »}
donc est pl.\ fidèle. Ainsi $R$ est connu \add{à équiv.\ près} quand on
connaît \uncertain{son image} \add{dans} $\widehat{S^{\circ}}$ \ill{}
\struck{\ill{}}.

Pour résumer, si $T$ est un $\lambda$-type ($\overrightarrow{\lambda}
= \emptyset$), alors $S = T(\mathrm{Ens})$ est muni d'une
sous-catégorie \struck{\ill{}} pleine $\Sigma \subset S$, formée des
$\varphi \in \mathrm{Ob}\,S$ tels que pour tout $C \in \mathrm{Ob}\,
\mathrm{Cat}_{\lambda}$ et tout \struck{\ill{}}
\uncertain{$\xi \in C_{\lambda}$} \note{on attendrait $\xi \in T(C)$},
\struck{\ill{}} \add{exprimant un} \struck{\ill{}} foncteur
\[
  \widetilde{\xi} : C^{\circ} \longrightarrow S, \qquad x \longmapsto
  \mathrm{Hom}(x, \xi)
\]
le foncteur composé avec $\mathrm{Hom}_{S}(\varphi, -)$ soit
représentable \note{un terme biffé dans le second argument} :
\[
  x \longmapsto \mathrm{Hom}_{S}(\varphi, \mathrm{Hom}(x, \xi))
\]
$\Sigma$ est stable dans $S$ par $\varinjlim$ \add{(inductives)} de
type $\in \lambda$ \note{« de type $\in \lambda$ » est peut-être
biffé} [car $C$ est stable dans $\hat{C}$ par $\varprojlim$ de type
$d$, $\forall C \in \mathrm{Ob}\,\mathrm{Cat}_{\lambda}$] et pour tt $C$
comme dessus, on a des foncteurs \note{au-dessus : « comme
$\lambda$-foncteurs »}
\[
  \Sigma^{\circ} \longrightarrow \underline{\mathrm{Hom}}(T(C), C),
  \qquad T(C) \longrightarrow \mathrm{Hom}_{\lambda}(\Sigma^{\circ}, C)
\]
On voit alors que l'on trouve ainsi une
\note{la phrase s'arrête là ; le reste de la page est blanc, et la
copie paginée ①--⑩ se termine ici}
\page{108}
\note{feuillet au crayon, d'une autre main de plume que la copie ①--⑩ ;
deux ajouts à l'encre bleue sont signalés ; les notations $B_{T}$,
$O_{B}$, $T(\mathrm{Ens})$ sont celles des pp. 97--106}
Soit $C$ une catégorie ess.\ petite avec $\varprojlim$ finies, $T$
l'espèce de structures que $C$ définit, $\hat{C} = B_{T} = B$ le topos
classifiant, $O_{B}$ la structure universelle dans $B$ (s'identifiant
à l'inclusion $C \to \hat{C}$). $E$ un topos.
\[
  O_{E} \in \mathrm{Ob}\,T(E) \simeq \mathrm{Ob}\,\mathrm{Hom}_{sex}(C, E)
  \ni f_{0}
\]
\note{sous $O_{E}$, une double flèche verticale $\Updownarrow$ menant
à la ligne suivante ; à droite, à l'encre bleue : « le morphisme
correspondant »}
\[
  f : E \longrightarrow B_{T}
\]
\note{devant $f$, un symbole raturé}
\struck{$\xi$ la structure} donc
\[
  O_{E} \simeq f^{*}_{T}(O_{B}) \qquad f_{0} = f^{*} \mid C
\]
$U \in \mathrm{Ob}\,E$, $X \in \mathrm{Ob}\,B$ (p.ex.\ $X \in
\mathrm{Ob}\,C$), on définit une \struck{bijection} \add{application}
\begin{equation}
  \mathrm{Hom}_{E}(U, f^{*}X) \longrightarrow
  \mathrm{Hom}(O_{B}(X), O_{E}(U))
  \tag{1}
\end{equation}
\note{la flèche porte un signe raturé ; sous le membre de droite, à
l'encre bleue : $T(\mathrm{Ens}) = \mathrm{Ind}(C^{\circ})$}
Où $O_{E}(U) \in T(\mathrm{Ens})$ est l'image de $O_{E} \in T(E)$ par le
foncteur $Z \mapsto \mathrm{Hom}_{E}(U, Z) : E \to (\mathrm{Ens})$ (qui
commute aux $\varprojlim$) et où on définit de $=$ $O_{B}(X)$ dans
$\hat{C}$. En effet, on a
\[
  \mathrm{Hom}_{E}(U, f^{*}(X)) \longrightarrow
  \mathrm{Hom}_{T(\mathrm{Ens})}(O_{E}(f^{*}(X)), O_{E}(U))
  \longrightarrow \mathrm{Hom}_{T(\mathrm{Ens})}(O_{B}(X), O_{E}(U))
\]
car on a un morphisme de $T(\mathrm{Ens})$
\[
  O_{B}(X) \longrightarrow O_{E}(f^{*}(X))
\]
\note{au bas de la page, un schéma sur deux lignes :}
\begin{gather*}
  C \longrightarrow \hat{C} \xrightarrow{\ h_{X}\ } (\mathrm{Ens}),
  \qquad Y \longmapsto \mathrm{Hom}_{C}(X, Y) =
  \mathrm{Hom}_{C^{\circ}}(Y, X) \\
  \to \mathrm{Hom}_{\hat{C}}(X, Y)
  \to \mathrm{Hom}_{E}(f^{*}(X), f^{*}(Y))
\end{gather*}
\note{au-dessous, biffé : $f_{0} \to E \xrightarrow{h_{f^{*}(X)}}$,
$Y \mapsto \mathrm{Hom}_{E}(\ldots, f_{0}(Y))$, le premier argument
raturé}

\page{109}
Prouver que si $X \in \mathrm{Ob}\,C$, l'application précédente est
une bijection. D'ailleurs, le 2ème membre s'exprime comme
$\mathrm{Hom}_{\mathrm{Ind}(C^{\circ})}(X^{\circ}, O_{E}(U))$ pour $U$
fixé ; la catégorie des couples $(X^{\circ}, \xi)$ ($\xi \in
\mathrm{Hom}(X^{\circ}, O_{E}(U))$) est la catégorie d'indices naturelle
définissant \uncertain{l'Ind-objet} $O_{E}(U)$ de
$\mathrm{Ind}(C^{\circ})$\ldots
\[
  O_{B}(X) : C \longrightarrow (\mathrm{Ens}), \qquad
  Y \longmapsto \mathrm{Hom}_{C}(X, Y) = \mathrm{Hom}_{C^{\circ}}(Y, X)
\]
\[
  O_{E}(U) : \qquad\qquad
  Y \longmapsto \mathrm{Hom}_{E}(U, f_{0}(Y))
\]
\note{le premier argument du dernier $\mathrm{Hom}$ est surchargé}
\[
  \mathrm{Hom}(O_{B}(X), O_{E}(U)) = O_{E}(U)(X) =
  \mathrm{Hom}_{E}(U, f_{0}(X))
\]

On va en profiter pour calculer $f^{*}(\mathcal{F})$, i.e.\
\uncertain{$\mathcal{F}$ un champ sur $B$}. Il revient au même de
\add{se} donner la restriction de $\mathcal{F}$ à \ill{} \ill{} $C$,
i.e.\ des données \ill{} une catégorie fibrée sur $C$. [\underline{NB}
si on a un sous-topos \add{$B'$} de $B = \hat{C}$, correspondant à une
topologie $\tau$ sur $C$, et à une sous-espèce de structure $T'$ de
$T$, alors si on \ill{} \uncertain{de cette espèce} i.e.\ \add{que}
$f$ \uncertain{se factorise} \add{par} $B'$ i.e.\ que $f_{0}$ est
continu \uncertain{pour} $\tau$, alors la donnée d'un champ
\add{$\mathcal{F}'$} sur $B'$ revient \struck{\ill{}} \add{à} la
donnée d'un champ $\mathcal{F}_{0}$ sur $B$ muni de \ill{} \ill{} de
$\tau$, et \ill{} \uncertain{notes} \ill{} \add{$f'^{*}$}$(\mathcal{F}')$
en termes de $f_{0}$ et $\mathcal{F}'$)] \struck{\ill{} que la
catégorie fibrée} On \struck{trouve que} \add{commence par le}
\uncertain{foncteur} \ill{} \note{à droite, sur la ligne : « \ill{}
un \uncertain{Foncteur} \uncertain{fibré} sur »} \struck{\ill{}}
$\mathcal{F}_{0}$ \add{le foncteur} \ill{} $\mathrm{Pro}(C) \simeq
T(\mathrm{Ens})^{\circ}$, \struck{\ill{}} \ill{} $\mathcal{F}_{0}$,
donc $\mathcal{F}_{0}(O_{E}(U))$ est défini [mais covariant en l'objet
$(\ )\in T(\mathrm{Ens})$]. Pour $U$ variable, $U \mapsto
\mathcal{F}_{0}(O_{E}(U))$ donne donc une catégorie fibrée sur $U$
[i\ldots\ image inverse de $\mathcal{F}_{0}$ par $O_{E} : E^{\circ} \to
T(\mathrm{Ens}) \cong \mathrm{Pro}(C)^{\circ}$] On prend le
\struck{catégorie} champ associé.
\note{page au crayon très pâle : la moitié des mots de liaison est
perdue, les formules passent}

\page{111}
\note{feuillet quadrillé à spirale, à l'encre ; titre souligné}
Topos modulaire pour des structures algébriques
\struck{\ill{}} sur objets \uncertain{variés} d'un topos fixé
\struck{\ill{}} $X$

On peut se donner \struck{\ill{}} des objets $E_{i} \in X$, $i \in I$,
et \struck{\ill{}} une espèce de structure $T$ avec topos modulaire $B$,
\struck{sur objets} avec données « objets de base » corresp.-dants : des
$b_{i} \in B$, $i \in I$. On s'intéresse aux, pour tt topos $X'$ sur
$X$, (aussi \uncertain{comme} images inverses des $E_{i}$ les $E'_{i}$)
\add{des} structures d'espèce $T$ \struck{dans} $X'$ \struck{que} (i.e.\
morphismes $f : X' \to B$) \struck{tels que} avec $f^{*}(b_{i}) \simeq
E'_{i}$. La catégorie de ces objets sur un topos $X'$ fixé, pour $X'
\to X$ et $X' \to B$ variables, est le produit 2-fibré de
$\underline{\mathrm{Hom}}_{\mathrm{Top}}(X', X)$ et
$\underline{\mathrm{Hom}}_{\mathrm{Top}}(X', B)$ sur
$\underline{\mathrm{Hom}}_{\mathrm{Top}}(X', B_{I})$. Donc on doit avoir
un \uncertain{produit} \uncertain{$B \times_{B_{I}} X$} dans la
2-catégorie $(\mathrm{Top})$

\begin{tikzcd}
  B \arrow[dr, "b"'] & & X \arrow[dl] \\
  & \widehat{(\mathrm{Ens}_{f})^{\circ}}{}^{I} = B_{I} &
\end{tikzcd}

\note{à droite du diagramme : « 2-fibre $B \times_{B_{I}} X$ »} (On
n'a pas encore \uncertain{prouvé} que cela marche.) \struck{\ill{}}
\struck{Cela doit donc \ill{}} supposons qu'il s'agisse plutôt
\add{$T$} d'une \struck{$\varprojlim$ finies} \add{structure définie
par $\varprojlim$ finies}, avec cat.\ modulaire $R$ (donc $B =
\hat{R}$, $I \to \mathrm{Ob}\,R$) Donc on s'intéresse aux foncteurs
exacts à g.\ $R \to X'$ transformant les $r_{i}$ en les $E'_{i}$. Soit
\struck{$\mathcal{S}$} \struck{la catégorie} \add{le champ sur $X$} des
\ill{} \struck{catégorie} formée des structures $(X', (E'_{i}))$
\add{et les foncteurs induits} \struck{i.e.\ des foncts ex\ldots}
\add{avec} $f : R \to X'$ \add{avec} $f(r_{i}) \simeq E'_{i}$ donnés, et
\struck{soit} \add{soit} $\underline{I}_{X}$ le champ \struck{sous la}
\add{la} \uncertain{fibre} \add{(en topos)} des \struck{$B$}.
\marginal{champ en cat.\ \ill{} (associé aux cat.\ \ill{})}
\note{note marginale en diagonale, au bas à gauche, barrée d'un
trait ; les insertions interlinéaires de ce feuillet sont nombreuses
et l'ordre de lecture est reconstitué}

\end{document}
